\subsubsection{Tensor Product}\label{subsubsec:tensor-product}

Let's start with something simple. As in previous sections, we have a simple one-qubit state. Now, we want to \textbf{add a second qubit to the system} to create a multi-qubit state. For simplicity, we will use only two qubits. \emph{How do we describe the \textbf{combined system}?} Simply keeping them separate is \underline{not} enough because we would miss correlations between the two qubits. \hl{We need a way to combine the two qubits into a \textbf{single object}.} The tool for doing so is the \textbf{tensor product}.

\highspace
\begin{examplebox}[: Dice Analogy]
    Imagine we have two dice. Each die can be in one of six states (1, 2, 3, 4, 5, or 6). If we want to describe the combined state of both dice, we can't just list their individual states separately. Instead, we need to consider all possible combinations of their states. The tensor product allows us to create a new space that includes all these combinations, giving us a complete description of the system.
\end{examplebox}

\highspace
The \definition{Tensor Product} is a \textbf{mathematical operation} that \hl{takes two vectors} (or more generally, two mathematical objects):
\begin{equation*}
    \ket{v_A} = a_0\ket{0} + a_1\ket{1} \quad \ket{v_B} = b_0\ket{0} + b_1\ket{1}
\end{equation*}
And \hl{produces a new vector} that represents the combined state of the two systems:
\begin{equation}
\begin{aligned}
    \ket{v_A} \otimes \ket{v_B} &= (a_0\ket{0} + a_1\ket{1}) \otimes (b_0\ket{0} + b_1\ket{1}) \\[.3em]
    &= a_0b_0\ket{00} + a_0b_1\ket{01} + a_1b_0\ket{10} + a_1b_1\ket{11}
\end{aligned}
\end{equation}
Here, the resulting state $\ket{v_A} \otimes \ket{v_B}$ is a \textbf{superposition of all possible combinations of the states of the two qubits}.

\highspace
\textcolor{Green3}{\faIcon{question-circle} \textbf{Why amplitudes multiply? For example, why is the amplitude of $\ket{00}$ equal to $a_0b_0$?}} A qubit is a vector:
\begin{equation*}
    \ket{v_{A}} = \begin{pmatrix}
        a_{0} \\
        a_{1}
    \end{pmatrix},
    \quad
    \ket{v_{B}} = \begin{pmatrix}
        b_{0} \\
        b_{1}
    \end{pmatrix}
\end{equation*}
The tensor product is \textbf{defined} as:
\begin{equation}
    \ket{v_{A}} \otimes \ket{v_{B}} = \begin{pmatrix}
        a_{0}b_{0} \\
        a_{0}b_{1} \\
        a_{1}b_{0} \\
        a_{1}b_{1}
    \end{pmatrix}
\end{equation}
So amplitudes multiply because \textbf{that's how the tensor product is defined}. \hl{\emph{But why is it defined like that?}} Because we want a rule that satisfies two requirements:
\begin{enumerate}
    \item \textbf{Linearity}. Quantum mechanics is linear:
    \begin{equation*}
        \left(a_{0} \ket{0} + a_{1} \ket{1}\right) \otimes \ket{v_{B}} = a_{0} \left(\ket{0} \otimes \ket{v_{B}}\right) + a_{1} \left(\ket{1} \otimes \ket{v_{B}}\right)
    \end{equation*}
    The tensor product must be \textbf{bilinear}, i.e., linear in each argument, so that superpositions are preserved when combining quantum states.
    \item \textbf{Basis consistency}:
    \begin{equation*}
        \ket{0} \otimes \ket{0} = \ket{00}, \quad \ket{0} \otimes \ket{1} = \ket{01}, \quad \ket{1} \otimes \ket{0} = \ket{10}, \quad \ket{1} \otimes \ket{1} = \ket{11}
    \end{equation*}
    The tensor product must map basis states of individual systems to a consistent basis of the joint system (e.g., $\ket{i} \otimes \ket{j} = \ket{ij}$), preserving the structure of the state space.
\end{enumerate}
This structure ensures that probabilities factorize correctly for independent systems.

\begin{flushleft}
    \textcolor{Red2}{\faIcon{exclamation-triangle} \textbf{Not all states can be factored}}
\end{flushleft}
Given two qubits, their joint state lives in a 4-dimensional vector space. However, \textbf{not all states} in this 4-dimensional space can be written as a tensor product of two single-qubit states. Precisely, if a multi-qubit state can be written as a \textbf{tensor product}, then it is \textbf{separable}, and we can recover the individual qubit states. But if it is \textbf{not a tensor product}, then it is \textbf{entangled}, and we cannot assign well-defined quantum states to each qubit individually.

\highspace
In general, exists a \textbf{necessary (and sufficient) condition} for being a tensor product \underline{for two qubits}. Given a 2-qubit state:
\begin{equation*}
    \ket{\psi} = \left(c_{0}, c_{1}, c_{2}, c_{3}\right)
\end{equation*}
Where $c_{0}, c_{1}, c_{2}, c_{3}$ are the amplitudes of the basis states $\ket{00}, \ket{01}, \ket{10}, \ket{11}$, respectively. The state \hl{$\ket{\psi}$ can be written as a tensor product of two single-qubit states if and only if}:
\begin{equation}\label{eq:separability-condition}
    c_{0}c_{3} = c_{1}c_{2}
\end{equation}
Equivalently, if we reshape the coefficients into a matrix
\begin{equation*}
    C =
    \begin{pmatrix}
        c_0 & c_1 \\
        c_2 & c_3
    \end{pmatrix}
\end{equation*}
The state is separable if and only if $\det(C) = 0$. More generally, for $n$-qubit states, there is \textbf{no single simple scalar condition} that characterizes separability. A state is separable \textbf{if and only if it can be written as a tensor product of $n$ single-qubit states}, which is a more complex condition to verify. So, the \hl{2-qubit case is simpler to analyze}, but for larger systems, determining whether a state is separable or entangled can be more challenging.

\highspace
Let's make this more concrete with two examples:
\begin{itemize}
    \item[\textcolor{Green3}{\faIcon{check}}] Let's start from a 4D vector:
    \begin{equation*}
        \ket{\psi} = \begin{pmatrix}
            \frac{1}{2} \\[.3em]
            \frac{1}{2} \\[.3em]
            \frac{1}{2} \\[.3em]
            \frac{1}{2}
        \end{pmatrix}
        = \frac{1}{2} \left( \ket{00} + \ket{01} + \ket{10} + \ket{11} \right)
    \end{equation*}
    We can check if it is a tensor product by verifying the condition in \autoref{eq:separability-condition}:
    \begin{equation*}
        c_{0}c_{3} = \frac{1}{2} \cdot \frac{1}{2} = \frac{1}{4}, \quad c_{1}c_{2} = \frac{1}{2} \cdot \frac{1}{2} = \frac{1}{4}
    \end{equation*}
    Since $c_{0}c_{3} = c_{1}c_{2}$, the state is a tensor product and thus separable.
    
    Now, we can find the two single-qubit states whose tensor product gives us $\ket{\psi}$. We try:
    \begin{equation*}
        \ket{v_{A}} = \begin{pmatrix}
            a_{0} \\
            a_{1}
        \end{pmatrix},
        \quad
        \ket{v_{B}} = \begin{pmatrix}
            b_{0} \\
            b_{1}
        \end{pmatrix}
    \end{equation*}
    We want:
    \begin{equation*}
        \left(a_0 b_0, a_0 b_1, a_1 b_0, a_1 b_1\right) = \left(\frac{1}{2}, \frac{1}{2}, \frac{1}{2}, \frac{1}{2}\right)
        \quad \Rightarrow \quad
        \begin{cases}
            a_0 b_0 = \frac{1}{2} \\[.3em]
            a_0 b_1 = \frac{1}{2} \\[.3em]
            a_1 b_0 = \frac{1}{2} \\[.3em]
            a_1 b_1 = \frac{1}{2}
        \end{cases}
    \end{equation*}
    Solving this system of equations, we find:
    \begin{equation*}
        a_0 = b_0 = \frac{1}{\sqrt{2}}, \quad a_1 = b_1 = \frac{1}{\sqrt{2}}
    \end{equation*}
    And this works because:
    \begin{equation*}
        \begin{cases}
            a_0 b_0 = \frac{1}{\sqrt{2}} \cdot \frac{1}{\sqrt{2}} = \frac{1}{2} \\[.3em]
            a_0 b_1 = \frac{1}{\sqrt{2}} \cdot \frac{1}{\sqrt{2}} = \frac{1}{2} \\[.3em]
            a_1 b_0 = \frac{1}{\sqrt{2}} \cdot \frac{1}{\sqrt{2}} = \frac{1}{2} \\[.3em]
            a_1 b_1 = \frac{1}{\sqrt{2}} \cdot \frac{1}{\sqrt{2}} = \frac{1}{2}
        \end{cases}
    \end{equation*}
    Therefore:
    \begin{equation*}
        \ket{\psi} = \left[ \frac{1}{\sqrt{2}} \left(\ket{0} + \ket{1}\right) \right] \otimes \left[ \frac{1}{\sqrt{2}} \left(\ket{0} + \ket{1}\right) \right]
    \end{equation*}
    The global 2-qubit state can be written as \textbf{two independent qubit states}, confirming that it is separable.
    \begin{equation*}
        \ket{\psi} = \begin{pmatrix}
            \frac{1}{2} \\[.3em]
            \frac{1}{2} \\[.3em]
            \frac{1}{2} \\[.3em]
            \frac{1}{2}
        \end{pmatrix}
        = \left[\frac{1}{\sqrt{2}} \left(\ket{0} + \ket{1}\right)\right] \otimes \left[\frac{1}{\sqrt{2}} \left(\ket{0} + \ket{1}\right)\right]
    \end{equation*}

    \item[\textcolor{Red2}{\faIcon{times}}] \phantomsection\label{ex:non-separable} Let's start from a different 4D vector:
    \begin{equation*}
        \ket{\psi} = \frac{1}{\sqrt{2}} \begin{pmatrix}
            1 \\
            0 \\
            0 \\
            1
        \end{pmatrix}
        = \frac{1}{\sqrt{2}} \left( \ket{00} + \ket{11} \right)
    \end{equation*}
    We check the condition for separability:
    \begin{equation*}
        c_{0}c_{3} = \frac{1}{2}, \quad c_{1}c_{2} = 0
    \end{equation*}
    Since $c_{0}c_{3} \neq c_{1}c_{2}$, the state is not a tensor product and thus is \textbf{not separable}. This state is an example of an \textbf{entangled state}, where the two qubits are correlated in such a way that we cannot describe them independently. We will see more about entanglement in the next section, but for now, the key takeaway is that not all multi-qubit states can be factored into tensor products of single-qubit states, and this is a fundamental aspect of quantum mechanics that leads to phenomena like entanglement.
\end{itemize}
In summary, the tensor product generates only a subset of the full 4-dimensional space of two-qubit states. Separable states have a special multiplicative structure and lie on a lower-dimensional manifold within this space. Most states in the full space do not satisfy this structure and are therefore entangled.